\section{Odd-$k$ and Even-$k$ Separations}
\label{sec:k-extension}

The two-point $1$-NN witness of \Cref{sec:finite-separation} shows that LOO
stability and replace-one stability are not interchangeable for $k=1$. A natural
question is whether the same phenomenon persists for larger $k$. We answer this
question in the affirmative: a general construction works for every odd $k \ge
3$, and an analogous (but tie-breaking-dependent) construction works for even
$k$.

\subsection{Odd-$k$ Separation}

Let $k = 2m+1$ for an integer $m \ge 1$. Let the metric space consist of two
distinct points $a$ and $b$ with $d(a,b) > 0$; any metric space with at least
two points suffices.

Define the ordered sample
\[
S_m \;=\; \big((a,0)_1,\; \ldots,\; (a,0)_{m+1},\;
          (a,1)_1,\; \ldots,\; (a,1)_{m},\;
          (b,0)_1,\; \ldots,\; (b,0)_{m}\big),
\]
where subscripts track the occurrence index within the ordered tuple.
The sample has size $n = (m+1) + m + m = 3m+1$.

Fix the query-label pair $(x,y) = (a,0)$ and the neighborhood size $k = 2m+1$.

\begin{theorem}[Odd-$k$ LOO/replace-one separation]
\label{thm:odd-k-separation}
For every odd $k = 2m+1$ with $m \ge 1$, the sample $S_m$ satisfies
\[
\Delta_{\mathrm{loo}}(S_m, i) = 0 \qquad \text{for every index } i,
\]
while there exists an index $j$ and a replacement point $z$ such that
\[
\Delta_{\mathrm{rep}}(S_m, j, z, a, 0) = 1.
\]
Thus, pointwise LOO stability does not imply pointwise replace-one stability
for deterministic $k$-NN with odd $k \ge 3$.
\end{theorem}

\begin{proof}
We analyze the margin behavior at the query $a$ and at each deleted occurrence.
Recall from \Cref{prop:5.1} that a deterministic $k$-NN prediction changes
if and only if the signed vote margin $M_k(S, x)$ crosses zero.

\paragraph{Original prediction at $(a,0)$.}
At query $a$, the ordered neighbors are sorted by distance then occurrence index.
All occurrences at $a$ are at distance $0$; all occurrences at $b$ are at
distance $d(a,b) > 0$. The deterministic order therefore selects all
$(m+1)$ copies of $(a,0)$ first, then all $m$ copies of $(a,1)$, giving exactly
$k = 2m+1$ neighbors. The signed margin is
\[
M_k(S_m, a) = (m+1)(-1) + m(+1) = -1 < 0,
\]
so the classifier predicts $0$ at $a$. Hence the loss at $(a,0)$ is $0$.

\paragraph{LOO at an $(a,0)$ occurrence.}
Delete the $i$-th copy of $(a,0)$ and evaluate at the deleted occurrence
$(a,0)$. After deletion the remaining sample has size $n-1 = 3m$, and
$k' = \min(k, n-1) = 2m+1$ for all $m \ge 1$.

At query $a$, the top-$k'$ neighbors consist of the remaining $m$ copies
of $(a,0)$, all $m$ copies of $(a,1)$, and the first copy of $(b,0)$.
The margin is
\[
M_{k'}(S_m^{-i}, a) = m(-1) + m(+1) + 1(-1) = -1.
\]
The prediction is $0$, unchanged from the original. Since the loss at $(a,0)$
was $0$ before and remains $0$ after deletion, $\Delta_{\mathrm{loo}} = 0$.

\paragraph{LOO at an $(a,1)$ occurrence.}
Delete the $i$-th copy of $(a,1)$ and evaluate at $(a,1)$. After deletion,
the top-$k'$ at $a$ consists of all $(m+1)$ copies of $(a,0)$, the remaining
$(m-1)$ copies of $(a,1)$, and the first copy of $(b,0)$. The margin is
\[
M_{k'}(S_m^{-i}, a) = (m+1)(-1) + (m-1)(+1) + 1(-1) = -3.
\]
The prediction is $0$. The original prediction at $(a,1)$ was also $0$
(margin $-1$), so the original loss at $(a,1)$ was $1$; the loss after deletion
is still $1$. Hence $\Delta_{\mathrm{loo}} = 0$.

\paragraph{LOO at a $(b,0)$ occurrence.}
Delete the $i$-th copy of $(b,0)$ and evaluate at $(b,0)$. The original margin
at query $b$ is
\[
M_k(S_m, b) = m(-1) + (m+1)(-1) = -(2m+1),
\]
so the original prediction is $0$ and the loss at $(b,0)$ is $0$.

After deletion, the top-$k'$ at $b$ consists of $(m-1)$ copies of $(b,0)$,
all $(m+1)$ copies of $(a,0)$, and the first copy of $(a,1)$. The margin is
\[
M_{k'}(S_m^{-i}, b) = (m-1)(-1) + (m+1)(-1) + 1(+1) = -(2m-1).
\]
For all $m \ge 0$, this is $\le 0$, so the prediction remains $0$ and the
loss stays $0$. Thus $\Delta_{\mathrm{loo}} = 0$.

Since every index falls into one of the three cases above, we have
$\Delta_{\mathrm{loo}}(S_m, i) = 0$ for all $i$.

\paragraph{Replace-one instability.}
Now replace the first occurrence $(a,0)_1$ with $(a,1)$. The perturbed
sample is
\[
S_m^{1 \leftarrow (a,1)} = \big((a,1),\; (a,0)_2,\ldots,(a,0)_{m+1},\;
                             (a,1)_1,\ldots,(a,1)_m,\;
                             (b,0)_1,\ldots,(b,0)_m\big).
\]

At query $a$, the top-$k$ neighbors are now: $(a,1)$ at index $1$, then
$(a,0)_2$ through $(a,0)_{m+1}$ ($m$ copies), then $(a,1)_1$ through
$(a,1)_m$ ($m$ copies). All $k = 2m+1$ neighbors are at distance $0$. The
signed margin becomes
\[
M_k(S_m^{1 \leftarrow (a,1)}, a) = 1(+1) + m(-1) + m(+1) = 1 > 0,
\]
so the classifier now predicts $1$ at $a$. The loss at $(a,0)$ changes from
$0$ to $1$, giving $\Delta_{\mathrm{rep}}(S_m, 1, (a,1), a, 0) = 1$.
\end{proof}

The construction uses only the fact that $a$ and $b$ are distinct points with
$d(a,b) > 0$; the separation therefore holds in any metric space with at least
two points. The ordered occurrence sequence is essential: placing all $(a,0)$
copies before $(a,1)$ copies ensures the original margin is $-1$. Conflicting
labels at $a$ are the mechanism by which a single replacement changes the label
composition of the top-$k$ neighborhood without changing its geometric
membership.

The $k=1$ case, excluded from \Cref{thm:odd-k-separation}, is already covered
by the two-point witness of \Cref{sec:finite-separation}. Taken together, the
LOO/replace-one separation holds for every odd $k$: $k=1$ by the two-point
witness and $k \ge 3$ by \Cref{thm:odd-k-separation}.

\subsection{Even-$k$ Separation}

For even $k = 2m$ with $m \ge 1$, the same metric space admits a separation
analogous to the odd case, but with a different margin profile.

Define the ordered sample
\[
S_m^{\text{even}} \;=\; \big((a,0)_1,\ldots,(a,0)_m,\;
                            (a,1)_1,\ldots,(a,1)_m,\;
                            (b,0)_1,\ldots,(b,0)_m\big),
\]
of size $n = 3m$. Fix $k = 2m$ and query-label pair $(a,0)$.

\begin{theorem}[Even-$k$ LOO/replace-one separation]
\label{thm:even-k-separation}
For every even $k = 2m$ with $m \ge 1$, the sample $S_m^{\text{even}}$
satisfies
\[
\Delta_{\mathrm{loo}}(S_m^{\text{even}}, i) = 0 \quad \text{for every index } i,
\]
while there exists an index $j$ and a replacement point $z$ such that
\[
\Delta_{\mathrm{rep}}(S_m^{\text{even}}, j, z, a, 0) = 1.
\]
\end{theorem}

\begin{proof}
\textbf{Original prediction at $(a,0)$.} At query $a$, the top-$k$ neighbors are
$m$ copies of $(a,0)$ and $m$ copies of $(a,1)$. The signed margin is
\[
M_k(S_m^{\text{even}}, a) = m(-1) + m(+1) = 0,
\]
and by the deterministic tie-breaking rule ($0 \prec 1$) the prediction is $0$,
so the loss at $(a,0)$ is $0$.

\textbf{LOO at an $(a,0)$ occurrence.} After deletion, the top-$k'$ at $a$
consists of $(m-1)$ copies of $(a,0)$, $m$ copies of $(a,1)$, and the first
copy of $(b,0)$. The margin is
\[
M_{k'} = (m-1)(-1) + m(+1) + 1(-1) = 0,
\]
the prediction remains $0$ by tie-breaking, and the loss at $(a,0)$ is
unchanged. Hence $\Delta_{\mathrm{loo}} = 0$.

\textbf{LOO at an $(a,1)$ occurrence.} After deletion, the top-$k'$ at $a$
consists of $m$ copies of $(a,0)$, $(m-1)$ copies of $(a,1)$, and the first
copy of $(b,0)$. The margin is
\[
M_{k'} = m(-1) + (m-1)(+1) + 1(-1) = -2,
\]
the prediction is $0$. The original loss at $(a,1)$ was $1$; after deletion the
loss remains $1$. So $\Delta_{\mathrm{loo}} = 0$.

\textbf{LOO at a $(b,0)$ occurrence.} After deletion, the top-$k'$ at $b$
consists of $(m-1)$ copies of $(b,0)$, $m$ copies of $(a,0)$, and the first
copy of $(a,1)$. The margin is
\[
M_{k'} = (m-1)(-1) + m(-1) + 1(+1) = -(2m-2),
\]
which is $\le 0$ for all $m \ge 1$. The loss at $(b,0)$ is unchanged.
Therefore $\Delta_{\mathrm{loo}} = 0$.

\textbf{Replace-one instability.} Replace the first occurrence $(a,0)_1$ with
$(a,1)$. The top-$k$ at $a$ now contains one copy of $(a,1)$ at the top,
$(m-1)$ copies of $(a,0)$, and $m$ copies of $(a,1)$. The margin is
\[
M_k = 1(+1) + (m-1)(-1) + m(+1) = 2 > 0,
\]
so the prediction flips to $1$. The loss at $(a,0)$ goes from $0$ to $1$,
giving $\Delta_{\mathrm{rep}} = 1$.
\end{proof}

\subsection{Structural Differences Between Odd and Even $k$}

The odd and even constructions share the same metric space and the same
replacement mechanism, but the margin analysis reveals a structural difference.

\begin{enumerate}[leftmargin=*,label=(\arabic*)]

\item \textbf{Margin at the prediction threshold.}
For odd $k$, the original margin is $-1$, strictly negative. The prediction
reflects a genuine (if minimal) majority. For even $k$, the original margin is
$0$, an exact tie; the prediction is determined solely by the tie-breaking
convention. In this sense, the even-$k$ separation is weaker: it depends on the
deterministic tie-breaking rule to stabilize the original prediction against
LOO perturbations.

\item \textbf{LOO stability margin.}
For odd $k$, after deleting an $(a,0)$ occurrence, the margin remains $-1$
(non-zero, no tie). For even $k$, the post-deletion margin becomes $0$
(a tie), and LOO stability relies on the same tie-breaking rule to keep the
prediction at $0$. If the tie-breaking rule were randomized, the even-$k$
prediction after deletion would be $0$ or $1$ with equal probability,
potentially breaking the LOO stability guarantee.

\item \textbf{Replace-one effect.}
In both cases, the margin increases by exactly $2$ after the replacement
(from $-1$ to $+1$ for odd $k$; from $0$ to $+2$ for even $k$). The increment
of $2$ follows from flipping one label from $0$ to $1$ in the top-$k$: the
contribution of that occurrence changes from $-1$ to $+1$.

\end{enumerate}

\subsection{Computational Evidence}

The constructions above subsume the computational search evidence for
$k=3,5,7$ that motivated this investigation. Those candidates were generated
via bounded enumeration and are now superseded by the uniform combinatorial
argument that extends the separation to all odd $k \ge 3$.

A summary of the computational candidates is provided in
\Cref{tab:k-gadget-candidates} in the appendix.
